\section{Research Methods}
% =============================================================================
% RESEARCH METHODS - v3 EMERGENT SOCIAL STRUCTURES FRAMING
% =============================================================================
% STATUS: Needs update to reflect emergent structures focus
%
% KEY CHANGES FROM v2:
%   1. Agents optimize ABSOLUTE resources (not relative position)
%   2. Game designed to enable RICH emergent structures (not just security dilemma)
%   3. Metrics expanded for coalition/hierarchy/norms (not just cooperation/conflict)
%   4. Two-comparison design still valid (architecture vs framing effects)
%
% STRUCTURE:
%   1. Method Justification - why this game enables emergent structures
%   2. Experimental Design - two-comparison design (keep from v2)
%   3. Computational Model Design - emphasize what each action enables
%   4. Agent Architectures - all optimize absolute resources
%   5. Measurements - expanded for full emergent structure analysis
%   6. Validation - frame around mechanism comparison, not theory validation
% =============================================================================

To answer the

\subsection{Method Justification}
% =============================================================================
% JUSTIFICATION - v3 FRAMING
% =============================================================================
% KEY POINTS:
%   1. Game designed to enable rich emergent structures
%      - Not testing one specific phenomenon (security dilemma)
%      - Testing what KINDS of structures emerge from different reasoning
%   2. Why RL: Pure optimization, trial-and-error learning
%      - Represents what Axelrod studied - mechanistic cooperation
%   3. Why LLM: Semantic reasoning with priors
%      - Norm recognition, identity formation, narrative interpretation
%      - Tests whether these enable DIFFERENT structures
%   4. General Sufficiency: minimal game, maximal emergence potential
%
% EXAMPLE DIRECTION:
%   "The game is designed following the General Sufficiency principle:
%   minimal elements that enable maximal emergent complexity. Each action
%   serves a distinct function: Invest-Self enables accumulation, Invest-Other
%   enables trust signaling, Arm-Self enables threat/defense, Arm-Other
%   enables coalition formation, Attack enables conflict and hierarchy.
%
%   This design allows the AGENTS to determine what structures emerge,
%   rather than the game mechanics forcing particular outcomes."
% =============================================================================

%Model justification
General Sufficiency principle -> as minimal as possible.
How do we produce, with as little elements as possible, the intended macro-behviour that we want to investigate -> cooperation, coalition, power hierarchies,
How do I make sure it is not void of meaning. Merely driven by my intuitions as someone that doesnt even study IR?

How do I make the model theoretically grounded, not just my ad hoc intuitions?

% TODO: Rewrite around emergent structures framing:
%   - Game enables multiple possible structures (not just one phenomenon)
%   - Each action serves distinct function for emergence
%   - Minimal design, maximal emergence potential
%   - Agents determine structures, not game mechanics

\subsubsection{Computational Mapping}
% =============================================================================
% REFRAME: Not about theories, about reasoning architectures
% =============================================================================
% RL = Optimization / Trial-and-error
%   - What Axelrod studied
%   - Discovers cooperation mechanistically
%   - No access to "meaning" of actions
%
% LLM = Semantic reasoning
%   - NEW capability: semantic priors
%   - Norm recognition, identity, narrative interpretation
%   - Question: does this produce different structures?
% =============================================================================

RL is a widely used technique for modelling IR theories, specifically useful for materialist views such as neorealism (as we can optimize our own resources).
Why?

% TODO: Rewrite to focus on reasoning architectures:
%   - RL: trial-and-error optimization (Axelrod tradition)
%   - LLM: semantic reasoning with priors (new capability)
%   - Not about "representing theories" but about comparing reasoning processes

CONSTRUCTIVISM
Core thesis
- identity formed through interaction, norms and group identity
- interests flow from identity

LLM justification
- LLMs have shown to adhere to norms and group identity through interaction (norms paper,)
- LLMs show different behaviour through contextual framing (personas, language!)

% TODO: Reframe as "semantic priors" not "constructivism":
%   - LLMs bring norm recognition, identity formation, narrative interpretation
%   - These are CAPABILITIES, not theory commitments
%   - LLM-Control uses these capabilities for materialist reasoning
%   - LLM-Constructivist uses them for identity/norm reasoning


%assumptions

\subsection{Experimental Design}
% =============================================================================
% TWO-COMPARISON DESIGN (keep from v2, still valid)
% =============================================================================
% COMPARISON 1: Architecture Effect (RL vs LLM-Control)
%   - SAME goal: maximize absolute resources
%   - DIFFERENT: how they reason
%   - Tests: does semantic reasoning ITSELF produce different structures?
%
% COMPARISON 2: Framing Effect (LLM-Control vs LLM-Constructivist)
%   - SAME architecture: both LLM
%   - DIFFERENT: how goals are framed (materialist vs identity/norms)
%   - Tests: do identity/norm priors add something beyond semantic reasoning?
%
% KEY UPDATE: All agents optimize ABSOLUTE resources (not relative)
%   - This makes cooperation naturally attractive
%   - Tests whether agents DISCOVER cooperation vs being forced away from it
%   - Divergence becomes more meaningful
% =============================================================================

% TODO: Write this section explaining:
%   1. The confounding problem (architecture vs goals)
%   2. Two-comparison solution
%   3. KEY: All agents optimize ABSOLUTE resources
%   4. Interpretation table for all four outcome combinations

\subsection{Computational Model Design}
% =============================================================================
% GAME DESIGN - v3: Emphasize what each action ENABLES
% =============================================================================
% The game is designed to enable rich emergent structures:
%
% | Action      | Effect                   | Enables...                    |
% |-------------|--------------------------|-------------------------------|
% | Invest-Self | Resource growth          | Individual accumulation       |
% | Invest-Other| Transfer resources       | Trust signaling, cooperation  |
% | Arm-Self    | Military capability      | Threat, defense, ambiguity    |
% | Arm-Other   | Alliance formation       | Coalitions, commitment        |
% | Attack      | Resource appropriation   | Conflict, hierarchy, dominance|
%
% POSSIBLE EMERGENT STRUCTURES:
%   - Bilateral cooperation (mutual Invest-Other)
%   - Multilateral coalitions (Arm-Other networks)
%   - Coalition wars (inter-group Attack)
%   - Free-riding (benefit from alliance without Arm-Other)
%   - Hierarchy (dominant agents extracting tribute)
%   - Norms (implicit rules like "don't attack investors")
%   - Polarization (system splits into opposing blocs)
% =============================================================================

The \textbf{State Space} consists of an array containing the resources of every agent and the history of all recent interactions, conflicts and their outcomes. Each round, all agents choose simultaneously from the \textbf{Action Space}:

\begin{table}[h]
\centering
\begin{tabular}{@{}llp{10cm}@{}}
\toprule
\textbf{Action} & \textbf{Target} & \textbf{Effect} \\
\midrule
Invest & Self & Costs $I$, generates $V$ (net positive growth) \\
Invest & Other & Costs $I$, other receives $V$ (signals trust) \\
\midrule
Arm & Self & Costs $D$, multiplier $M/O$ for 1 round (general threat) \\
Arm & Other & Costs $D$, multiplier $M/O$ for 2 rounds (signals alliance) \\
\midrule
Attack & Other & Probabilistic conflict resolution (see below) \\
\bottomrule
\end{tabular}
\end{table}

% TODO: Add paragraph explaining what each action ENABLES:
%   "Each action serves a distinct function in enabling emergent structures.
%   Invest-Self enables accumulation. Invest-Other enables trust signaling
%   and bilateral cooperation. Arm-Self creates ambiguity (defense or threat?).
%   Arm-Other enables coalition formation and commitment signals. Attack
%   enables conflict, hierarchy formation, and resource redistribution.
%
%   This design allows agents to organically develop cooperation, coalitions,
%   hierarchies, or conflict - the AGENTS determine what emerges, not the
%   game mechanics."

\noindent\textbf{Conflict Resolution:} The resources of both sides, consisting of the main combatants plus active Arms multipliers, are summed and relative to the total resources involved, win probabilties are determined.
The winner recieves $T\%$ of the losers resources. All combatants pay $C$ conflict cost.

% TODO: Add game structure:
%   "Agents play for N rounds with complete information about resource
%   holdings and recent action history (last K rounds). The game length
%   is sufficient to allow emergent structures to form and stabilize."

\subsection{Agent Architectures}
% =============================================================================
% AGENT ARCHITECTURES - v3: All optimize ABSOLUTE resources
% =============================================================================
% KEY CHANGE: Not "relative gain" but "absolute resources"
%   - Cooperation becomes naturally attractive (positive-sum possible)
%   - Tests whether agents DISCOVER cooperation
%   - Divergence is about reasoning, not reward structure
%
% ORDER: RL → LLM-Control → LLM-Constructivist
%   - LLM-Control is PRIMARY comparison to RL (architecture effect)
%   - LLM-Constructivist is SECONDARY comparison to LLM-Control (framing effect)
% =============================================================================

The simulation is done using homogenous agent architectures for every agent.

\textbf{Reinforcement Learning (RL) Agents} represent state actors in Waltz's theory of Neorealism. At it's core, the agents should be self-interested and concerned with optimizing their relative gain within the system. Computationally, this identity is translated through optimizing the following function during training:
[insert RL optimization function]

% TODO: Rewrite RL description (v3):
%   "RL agents optimize absolute resource accumulation through trial-and-error
%   learning. They discover strategies that maximize their resources over time,
%   potentially including cooperation if it proves beneficial. However, they
%   have no access to semantic interpretation of other agents' actions - they
%   learn purely from reward signals.
%
%   This represents the Axelrod tradition: can optimizers discover cooperation
%   mechanistically?"

\textbf{Large-Language Model (LLM) Agents} represent state actors in Wendt's theory of Constructivism. Based on the intersubjective interpretation of shared history, LLM agents are prompted internally update their identities and interests, and choose their next action accordingly.

% TODO: Rewrite as LLM-Constructivist (v3):
%   "LLM-Constructivist agents are prompted to develop their identity through
%   interaction and let their interests flow from that identity. They interpret
%   other agents' actions through normative frames and identity narratives.
%   Their goal is still resource accumulation, but their reasoning process
%   incorporates norm recognition and identity formation.
%
%   KEY: Same goal as RL and LLM-Control, but with identity/norm framing.
%   This is Comparison 2: Framing Effect."

\textbf{Control LLM Agents} represent the state actors in Neorealistic theory and are prompted to optimize their relative gain through their actions. These agents are specifically introduced to effectively isolate the effects of the constructivist prompting.

% TODO: Rewrite LLM-Control (v3) - move BEFORE LLM-Constructivist:
%   "LLM-Control agents use the same LLM architecture but with minimal
%   goal framing: maximize your resources. They can semantically interpret
%   other agents' actions ('is this arming defensive or aggressive?') but
%   without explicit identity/norm prompting.
%
%   KEY: Same goal as RL, same architecture as LLM-Constructivist.
%   This is Comparison 1: Architecture Effect.
%   If RL ≠ LLM-Control, semantic reasoning itself matters."

\subsection{Measurements}
% =============================================================================
% MEASUREMENTS - v3: EXPANDED FOR EMERGENT STRUCTURES
% =============================================================================
% OLD (v1/v2): Coalition, Hierarchy, Cooperation/Conflict (generic)
% NEW (v3): Full emergent structure analysis
%
% COALITION DYNAMICS:
%   - Number of distinct coalitions (clustering analysis)
%   - Coalition size distribution
%   - Coalition stability (duration of Arm-Other relationships)
%   - Coalition formation speed (rounds until first stable coalition)
%   - Inter-coalition vs intra-coalition conflict rates
%
% COOPERATION PATTERNS:
%   - Investment frequency (self vs other)
%   - Reciprocity rate: P(Invest-Other | received Invest-Other)
%   - Trust persistence: do cooperative relationships survive conflict?
%   - Cooperation network density and clustering
%
% HIERARCHY/DOMINANCE:
%   - Gini coefficient of resources over time
%   - Emergence of "hub" agents (high degree centrality)
%   - Tributary relationships (asymmetric investment flows)
%   - Power concentration trajectory
%
% CONFLICT DYNAMICS:
%   - Attack frequency over time
%   - Betrayal rate: P(Attack | former ally)
%   - Conflict escalation patterns
%   - Arms race indicators (Arm frequency spirals)
%
% NORM EMERGENCE (qualitative + quantitative):
%   - Do agents develop implicit rules? (e.g., "don't attack investors")
%   - Punishment of norm violators
%   - In-group/out-group treatment differences
%   - (For LLMs: analyze reasoning traces for norm references)
%
% NETWORK STRUCTURE:
%   - Modularity (do distinct communities form?)
%   - Clustering coefficient
%   - Assortivity (do similar agents cluster?)
%   - Network evolution over time
% =============================================================================

To clearly analyse the emergent social structures, three specific elements are analysed, each with specific metrics to track.

For \textbf{Coalition}:
\begin{itemize}
    \item Coalition stability: average duration of Arm-other relationships
    \item Trade stability: average duration of Invest-other relationships
    \item Reciprocity rate: P(A arms B | B armed A last round)
    \item Betrayal rate: P(attack | previous ally)
\end{itemize}

For \textbf{Hierarchy}:
\begin{itemize}
    \item Gini coefficient of resources
    \item Network centrality of Arm-other relationships
\end{itemize}

For \textbf{Cooperation/Conflict}:
\begin{itemize}
    \item Attack frequency per round
    \item Invest frequency per round
    \item Total welfare
\end{itemize}

% TODO: Expand measurements for v3 emergent structures analysis.
% Current metrics are good foundation but add:
%
% COALITION DYNAMICS (expand):
%   - Number of distinct coalitions
%   - Coalition formation speed
%   - Inter vs intra-coalition conflict
%
% NORM EMERGENCE (new):
%   - Implicit rule development
%   - In-group/out-group treatment differences
%   - For LLMs: analyze reasoning traces
%
% NETWORK STRUCTURE (new):
%   - Modularity
%   - Network evolution over time
%
% COMPARISON METRICS:
%   - Same metrics for all three conditions
%   - Statistical comparison: RL vs LLM-Control (architecture)
%   - Statistical comparison: LLM-Control vs LLM-Constructivist (framing)

\subsection{Validation}
% =============================================================================
% VALIDATION - v3 FRAMING
% =============================================================================
% NOT VALIDATING:
%   - Whether LLMs "realistically" simulate humans/states
%   - Micro-level behavioral accuracy
%   - Prediction of real-world outcomes
%
% VALIDATING:
%   1. Do different architectures produce different emergent structures?
%   2. Can we identify WHAT about semantic reasoning produces differences?
%      (Analyze LLM reasoning traces: norms? identity? trust narratives?)
%   3. Are results robust across parameter variations?
%
% FRAMING: Computational experiment testing whether semantic priors
% enable different social structures, not validating agents or theories.
% =============================================================================

The use of Generative Agents within ABMs has been criticed for enlarging problems around validation and calibration that are known in ABMs. \cite{larooij2025}. The thesis should make clear what is validated, and, importantly, what is not to prevent it being "void of meaning".

% TODO: Rewrite validation for v3:
%
% "Following Larooij & Törnberg (2025), this study does not claim that
% agents realistically simulate real actors. Instead, validation focuses on:
%
% 1. STRUCTURE COMPARISON: Do different agent architectures produce
%    measurably different emergent structures under identical game rules?
%
% 2. MECHANISM IDENTIFICATION: If structures differ, WHY? We analyze
%    LLM reasoning traces to identify whether norm references, identity
%    formation, or trust narratives explain the difference.
%
% 3. ARCHITECTURE ISOLATION: Does the RL vs LLM-Control comparison
%    show that semantic reasoning itself (not just prompting) matters?
%
% 4. ROBUSTNESS: Are patterns stable across parameter variations
%    (number of agents, game length, resource parameters)?
%
% This is a computational experiment asking: Do semantic priors enable
% different emergent structures? Not: Do these agents accurately model
% real-world actors?"
